\documentclass[11pt,a4paper]{article}
\usepackage[utf8]{inputenc}
\usepackage[ngerman]{babel}
\usepackage{geometry}
\geometry{left=2cm,right=2cm,top=2cm,bottom=2cm}
\usepackage{enumitem}
\usepackage{hyperref}
\usepackage{booktabs}
\usepackage{xcolor}

\pagestyle{empty}
\setlength{\parindent}{0pt}
\setlength{\parskip}{0.8ex}

\newcommand{\sectiontitle}[1]{
    {\large \textbf{#1}} \vspace{0.2cm} \hrule \vspace{0.5cm}
}

\begin{document}

% Header
{\Large \textbf{Victor Minig}} \\
\vspace{0.2cm}
\small
Geboren: 08.04.1998 \\
Wörthstraße 1c, 96052 Würzburg, Deutschland \\
+49 151 676 11588 \,|\, \href{mailto:victor.minig@web.de}{victor.minig@web.de} \\
\href{https://www.linkedin.com/in/victor-minig}{LinkedIn} \,|\, \href{https://github.com/}{GitHub} \,|\, \href{https://rickcler.github.io/VictorMinig.github.io/}{Portfolio-Website}

\vspace{1cm}

% % Kurzprofil
% \sectiontitle{Kurzprofil}
% Datenanalytiker mit sozial- und politikwissenschaftlichem Hintergrund (Abschlussnote 1,3) und Schwerpunkt auf quantitativen Methoden, statistischer Modellierung und Datenanalyse. Erfahrung in der Arbeit mit Paneldaten, Mehrebenenmodellen und Faktorenanalyse. Sehr gute Kenntnisse in Python (Pandas, NumPy, Seaborn), Stata und R. Selbstständige Arbeitsweise, hohe Lernbereitschaft und Interesse an statistischer Methodenentwicklung für Promotionsvorhaben.

\vspace{0.5cm}

% Berufserfahrung
\sectiontitle{Berufserfahrung}

\textbf{Wissenschaftliche Hilfskraft (HiWi)} \hfill \textbf{seit 11/2025} \\
Lehrstuhl für Kognitive Systeme, Universität Würzburg \\
\vspace{0.1cm}
\begin{itemize}[leftmargin=*, noitemsep, topsep=0pt]
\item Entwicklung und Implementierung eines Clustering-Algorithmus für hochdimensionale Prozessdaten zur Mustererkennung
\item Implementierung eines rekurrenten neuronalen Netzes (RNN) zur Analyse sequenzieller Prozessdaten
\end{itemize}
\vspace{0.3cm}

\textbf{Datenanalyse - Werkstudent} \hfill \textbf{02/2024 – 10/2025} \\
Siemens AG, Nürnberg  \\
\vspace{0.1cm}
\begin{itemize}[leftmargin=*, noitemsep, topsep=0pt]
\item Entwicklung eines Webscraping-Algorithmus in Python (Selenium) für den Bundesklinikatlas
\item Entwicklung von ETL-Prozessen in Python und SQL
\item Zusammenarbeit mit Fachabteilungen zur datengestützten Entscheidungsfindung
\end{itemize}
\vspace{0.3cm}

\textbf{Supply Chain Analyst} \hfill \textbf{05/2023 – 08/2023} \\
Flaschenpost SE, Münster \\
\vspace{0.1cm}
\begin{itemize}[leftmargin=*, noitemsep, topsep=0pt]
\item Entwicklung eines Reportings zur Generierung von Bestellvorschlägen für die interne Disposition unter Berücksichtigung verschiedener Stakeholder-Anforderungen
\item Erstellung eines Dashboards und Durchführung von Analysen zur Identifikation mangelhafter Artikelanzahlen bei Anlieferung
\item Entwicklung von Datenbankabfragen sowie Datenvisualisierung mit Power BI
\end{itemize}
\vspace{0.3cm}

\textbf{Praktikum Datenanalyst} \hfill \textbf{10/2021 – 02/2022} \\
Institut für Hochschulforschung, München \\
\vspace{0.1cm}
\begin{itemize}[leftmargin=*, noitemsep, topsep=0pt]
\item Analyse quantitativer Daten, deskriptive und bivariate Analysen demografischer Daten in Excel und Stata
\item Explorative Faktorenanalyse zur Aufdeckung latenter Datenstrukturen
\item Datenmanagement, -edition und Präsentation der Ergebnisse vor der Institutsleitung
\end{itemize}

\vspace{0.5cm}

% Ausbildung
\sectiontitle{Ausbildung}

\textbf{Master Survey Statistics and Data Analysis} \hfill \textbf{10/2023 – Heute} \\
Otto-Friedrich-Universität Bamberg
\vspace{0.3cm}

\textbf{Bachelor Political and Social Studies} \hfill \textbf{10/2018 – 09/2022} \\
Julius-Maximilians-Universität Würzburg, Abschlussnote: 1,3 \\
Vertiefung: Datenanalyse \\
Bachelorarbeit: \textit{Multilevel-Untersuchung der Determinanten von Lebenszufriedenheit in europäischen Ländern} (Note: 1,3)
\vspace{0.3cm}

\textbf{Auslandssemester} \hfill \textbf{10/2020 – 03/2021} \\
Universität Ljubljana, Slowenien
\vspace{0.3cm}

\textbf{Studium Mathematik (ohne Abschluss)} \hfill \textbf{10/2017 – 10/2018} \\
Julius-Maximilians-Universität Würzburg
\vspace{0.3cm}

\textbf{Work and Travel} \hfill \textbf{11/2016 – 09/2017}
\vspace{0.3cm}

\textbf{Abitur} \hfill \textbf{10/2016} \\
Matthias-Grünewald-Gymnasium, Würzburg

\vspace{0.5cm}

% Methoden & Statistik-Fokus
\sectiontitle{Statistische und methodische Kenntnisse}

\begin{itemize}[leftmargin=*, noitemsep, topsep=0pt]
\item \textbf{Programmierung / Software:} Python (Pandas, NumPy, Scikit-learn), R (tidyverse, lme4), Stata, Excel
\end{itemize}

\vspace{0.5cm}

% Sprachen
\sectiontitle{Sprachen}

\begin{tabular}{@{}ll}
Deutsch & Muttersprache \\
Englisch & Fließend (C1)
\end{tabular}

\end{document}